% problem-000381 9 化学反应与结构综合分析
\section*{9. 化学反应与结构综合分析}
$
\mathrm{A} + \mathrm{CH} _ {3} \mathrm{MgCl} \longrightarrow \xrightarrow {\mathrm{H} _ {2} \mathrm{O}} \quad \mathrm{M} (\mathrm{C} _ {8} \mathrm{H} _ {1 2} \mathrm{O}) \xrightarrow [ \Delta ]{\mathrm{H} _ {2} \mathrm{SO} _ {4}} \quad \mathrm{N} (\mathrm{C} _ {8} \mathrm{H} _ {1 0})
$

⽣成物N分⼦中只有3种不同化学环境的氢，它们的数⽬⽐为 $1 { : } 1 { : } 3 _ { \circ }$

\subsection*{9-1}
画出化合物A、B、C、D、E、M和N的结构简式。

\subsection*{9-2}
A、B和C互为哪种异构体？（在正确选项的标号前打钩）① 碳架异构体 ② 位置异构体 ③ 官能团异构体 ④ 顺反异构体

\subsection*{9-3}
A能⾃发转化为B和C，为什么？

\subsection*{9-4}
B和C在室温下反应可得到⼀组旋光异构体L，每个旋光异构体中有 个不对称碳原⼦。
